\documentclass[10pt,landscape,a4paper]{article}
\usepackage[utf8]{inputenc}
\usepackage[T1]{fontenc}
\usepackage{geometry}
\usepackage{multicol}
\usepackage{amsmath,amssymb}
\usepackage{enumitem}
\usepackage{graphicx}
\usepackage{xcolor}
\usepackage{listings}
\usepackage{booktabs}
\usepackage{titlesec}
\usepackage{fancyhdr}
\usepackage{hyperref}
\usepackage{CJKutf8}

% Page layout
\geometry{top=0.8cm, bottom=0.8cm, left=0.8cm, right=0.8cm}
\pagestyle{empty}
\setlength{\parindent}{0pt}
\setlength{\parskip}{2pt}
\setlength{\columnsep}{8pt}

% Section formatting
\titleformat{\section}{\normalsize\bfseries\color{blue!70!black}}{}{0em}{}[\titlerule]
\titleformat{\subsection}{\small\bfseries\color{blue!50!black}}{}{0em}{}
\titlespacing{\section}{0pt}{4pt}{2pt}
\titlespacing{\subsection}{0pt}{3pt}{1pt}

% Code listing style
\lstset{
    language=Python,
    basicstyle=\ttfamily\scriptsize,
    keywordstyle=\color{blue!70!black}\bfseries,
    stringstyle=\color{red!60!black},
    commentstyle=\color{green!50!black}\itshape,
    numberstyle=\tiny\color{gray},
    backgroundcolor=\color{gray!8},
    frame=single,
    framerule=0.3pt,
    rulecolor=\color{gray!50},
    breaklines=true,
    breakatwhitespace=false,
    tabsize=2,
    showstringspaces=false,
    aboveskip=3pt,
    belowskip=3pt,
    xleftmargin=2pt,
    xrightmargin=2pt,
}

% Custom box for key concepts
\newcommand{\keybox}[1]{%
    \begin{center}
    \fcolorbox{blue!50!black}{blue!5}{%
        \parbox{0.95\linewidth}{\small #1}%
    }
    \end{center}
}

% Chinese hint style (gray, smaller)
\newcommand{\zhint}[1]{{\scriptsize\color{gray!70!black}#1}}

\begin{document}
\begin{CJK}{UTF8}{gbsn}

\begin{center}
    {\large\bfseries IERG4320/IEMS5726 --- Chapter 1: Course Introduction --- Cheat Sheet}\\[2pt]
    {\small Prof.\ YAN Wenjing, Dr.\ Danny Yip \quad|\quad Department of Information Engineering, CUHK}
\end{center}

\vspace{2pt}
\begin{multicols}{3}

% ============================================================
\section{Course Overview \zhint{(课程概览)}}
% ============================================================

\begin{itemize}[leftmargin=*, nosep, itemsep=1pt]
    \item \textbf{Lecture}: Wed 7:00--10:00 pm @ LSB LT6 (12 lectures)
    \item \textbf{Final Exam}: 15 Apr, ERB 1004, open book/note, close internet
\end{itemize}

\subsection{Assessment \zhint{(评分占比)}}
\begin{tabular}{ll}
UReply attendance \zhint{(签到)} & 5\% \\
4 Assignments \zhint{(作业)} & 35\% \\
Course project \zhint{(项目, 截止5/8)} & 20\% \\
Final exam \zhint{(期末, 开卷)} & 40\% \\
\end{tabular}

\subsection{Fail Conditions \zhint{(挂科条件)}}
\begin{itemize}[leftmargin=*, nosep, itemsep=0pt]
    \item Final exam below 30\% \zhint{(期末低于30分)}
    \item Grand total below 40\% \zhint{(总分低于40分)}
    \item Not submitting course project \zhint{(不交项目直接挂)}
\end{itemize}

\subsection{AI Tool Policy \zhint{(AI工具政策)}}
ChatGPT/DeepSeek/Gemini \textbf{allowed} for coursework, but \textbf{NOT} for final exam.\\
\zhint{作业可以用AI，但必须注明哪些部分用了AI；期末考试不能用。}

% ============================================================
\section{Key Terminology \zhint{(核心术语)}}
% ============================================================

\subsection{Data-related Terms \zhint{(数据相关)}}
\begin{itemize}[leftmargin=*, nosep, itemsep=1pt]
    \item \textbf{Database} \zhint{(数据库)}: organized data collection\\
    \zhint{就是把数据整整齐齐存起来，方便查找}
    \item \textbf{Data Analysis} \zhint{(数据分析)}: get knowledge from data\\
    \zhint{从数据里挖出有用的信息}
    \item \textbf{Data Mining} \zhint{(数据挖掘)}: find patterns in \emph{large} datasets\\
    \zhint{在海量数据中找隐藏的规律}
    \item \textbf{Big Data} \zhint{(大数据)}: dealing with huge, complex data\\
    \zhint{数据量大到普通方法处理不了，需要专门技术}
\end{itemize}

\subsection{AI / ML / DL \zhint{(三者关系)}}
\begin{itemize}[leftmargin=*, nosep, itemsep=1pt]
    \item \textbf{AI} \zhint{(人工智能)}: machines that can think/act smart\\
    \zhint{让机器表现出类似人类的智能行为}
    \item \textbf{ML} \zhint{(机器学习)}: machines learn from data, no need to hard-code rules\\
    \zhint{不用写死规则，让机器自己从数据中学习；是AI的子集}
    \item \textbf{DL} \zhint{(深度学习)}: ML using deep neural networks (many layers)\\
    \zhint{用多层神经网络来学习；是ML的子集}
\end{itemize}

\keybox{
\textbf{Relationship}: AI $\supset$ ML $\supset$ DL\\
\zhint{AI最大，ML在AI里面，DL在ML里面，像套娃一样}
}

\subsection{Roles \zhint{(数据科学三大岗位)}}
\begin{itemize}[leftmargin=*, nosep, itemsep=1pt]
    \item \textbf{Data Scientist} \zhint{(数据科学家)}: turns data into insights \& decisions\\
    \zhint{从数据中找出有价值的结论，帮助决策}
    \item \textbf{ML Engineer} \zhint{(机器学习工程师)}: builds ML models into real apps\\
    \zhint{把模型落地，做成能用的产品}
    \item \textbf{Data Engineer} \zhint{(数据工程师)}: builds data pipelines \& infrastructure\\
    \zhint{搭数据管道，保证数据能顺畅地收集、存储、流转}
\end{itemize}

% ============================================================
\section{Data Science \zhint{(数据科学)}}
% ============================================================

\zhint{用科学方法从数据中提取知识和洞察，并应用到实际场景中。}

\subsection{Data Science Lifecycle \zhint{(生命周期)}}
\begin{enumerate}[leftmargin=*, nosep, itemsep=1pt]
    \item \textbf{Data Acquisition} \zhint{(采集)} --- collect raw data
    \item \textbf{Data Cleaning} \zhint{(清洗)} --- fix errors, remove noise
    \item \textbf{Data Modeling} \zhint{(建模)} --- build models to find patterns
    \item \textbf{Data Analysis} \zhint{(分析)} --- draw conclusions
    \item \textbf{Data Visualization} \zhint{(可视化)} --- show results with charts
\end{enumerate}

\keybox{
\zhint{简记：采集 $\to$ 清洗 $\to$ 建模 $\to$ 分析 $\to$ 可视化}
}

\subsection{Related Disciplines \zhint{(相关学科)}}
\begin{itemize}[leftmargin=*, nosep, itemsep=0pt]
    \item Computer Science \zhint{(计算机)} --- programming, databases
    \item Statistics \zhint{(统计学)} --- interpret results
    \item Domain knowledge \zhint{(领域知识)} --- depends on the problem
\end{itemize}

\subsection{Why Data Science Matters \zhint{(为啥重要)}}
\begin{enumerate}[leftmargin=*, nosep, itemsep=0pt]
    \item Cost efficient \zhint{(省钱)}
    \item Better decision making \zhint{(帮助做更好的决策)}
    \item Understand customers \zhint{(了解用户行为)}
    \item Manage risks \zhint{(风险管理)}
    \item Drive innovation \zhint{(推动创新)}
\end{enumerate}

% ============================================================
\section{Classic Examples \zhint{(经典案例)}}
% ============================================================

\begin{enumerate}[leftmargin=*, nosep, itemsep=1pt]
    \item \textbf{MNIST} \zhint{(手写数字识别)} --- one of the earliest ML systems\\
    \zhint{给机器看手写的0--9，让它自动认出是哪个数字}
    \item \textbf{Recommender Systems} \zhint{(推荐系统)} --- e.g.\ Netflix Prize\\
    \zhint{根据你的观影历史，猜你还想看什么电影}
    \item \textbf{Face Detection} \zhint{(人脸识别)} --- e.g.\ Microsoft Face API\\
    \zhint{从照片/视频里找到人脸并认出是谁}
    \item \textbf{Machine Translation} \zhint{(机器翻译)} --- e.g.\ Google Translate\\
    \zhint{自动把一种语言翻译成另一种}
    \item \textbf{Face Generator} \zhint{(人脸生成)} --- GANs\\
    \zhint{AI凭空生成逼真的人脸照片}
    \item \textbf{ChatGPT} \zhint{(聊天机器人)} --- NLP-powered chatbot\\
    \zhint{能像人一样对话，写邮件、写代码}
\end{enumerate}

% ============================================================
\section{Types of Data Analytics \zhint{(数据分析类型)}}
% ============================================================

\begin{enumerate}[leftmargin=*, nosep, itemsep=1pt]
    \item \textbf{Predictive} \zhint{(预测型)}: ``What will happen next?''\\
    \zhint{预测未来。例：明天要进多少巧克力冰淇淋？}
    \item \textbf{Descriptive} \zhint{(描述型)}: ``What happened?''\\
    \zhint{总结过去。例：上周最受欢迎的口味是什么？}
    \item \textbf{Prescriptive} \zhint{(建议型)}: ``What should we do?''\\
    \zhint{给建议。例：要不要提高巧克力冰淇淋的品质？}
    \item \textbf{Diagnostic} \zhint{(诊断型)}: ``Why did it happen?''\\
    \zhint{找原因。例：为什么巧克力和香草卖得最好？}
\end{enumerate}

% ============================================================
\section{AI Ethics \& Governance \zhint{(AI伦理与监管)}}
% ============================================================

\subsection{EU AI Act \zhint{(欧盟AI法案)}}
\zhint{全球第一部AI法规，按风险等级分类管理：}

\begin{center}
\scriptsize
\begin{tabular}{lp{4.2cm}}
\toprule
\textbf{Risk} & \textbf{Examples / Rules} \\
\midrule
Unacceptable \zhint{(禁止)} & Social scoring, deceptive AI $\to$ \textbf{banned} \\
High \zhint{(高风险)} & Medical AI, self-driving $\to$ strict review \\
Limited \zhint{(有限风险)} & ChatGPT, deepfakes $\to$ must be transparent \\
Minimal \zhint{(低风险)} & Spam filters, AI games $\to$ self-regulate \\
\bottomrule
\end{tabular}
\end{center}

\subsection{HK Ethical AI Framework \zhint{(香港AI伦理框架, 2024)}}
\zhint{数字政策办公室发布，12条原则：}

\begin{enumerate}[leftmargin=*, nosep, itemsep=0pt]
    \item Transparency \zhint{(透明)}
    \item Reliability \& Security \zhint{(可靠安全)}
    \item Fairness \zhint{(公平，不歧视)}
    \item Diversity \& Inclusion \zhint{(多元包容)}
    \item Human Oversight \zhint{(人类监督)}
    \item Lawfulness \zhint{(合法合规)}
    \item Safety \zhint{(安全)}
    \item Data Privacy \zhint{(数据隐私)}
    \item Accountability \zhint{(问责)}
    \item Beneficial AI \zhint{(有益于社会)}
    \item Cooperation \zhint{(开放合作)}
    \item Sustainability \zhint{(可持续发展)}
\end{enumerate}

\subsection{Mainland China \zhint{(中国大陆)}}
\zhint{2023年8月起实施《生成式AI服务管理暂行办法》，基于《网络安全法》《数据安全法》《个人信息保护法》。}

% ============================================================
\section{Course Approach \zhint{(课程特点)}}
% ============================================================

\zhint{这门课不只是普通的ML入门课，还会覆盖：}
\begin{itemize}[leftmargin=*, nosep, itemsep=0pt]
    \item Data collection \& preprocessing \zhint{(数据收集与预处理)}
    \item Data representation \zhint{(数据表示)}
    \item Data visualization \zhint{(数据可视化)}
    \item Domain-specific problems \zhint{(各领域实际问题)}
\end{itemize}

\textbf{Prerequisites} \zhint{(先修知识)}: Statistics, Linear Algebra, Calculus, Probability.\\
\zhint{需要统计学、线性代数、微积分、概率论基础。}

% ============================================================
\section{Course Schedule \zhint{(课程安排)}}
% ============================================================

\begin{center}
\scriptsize
\begin{tabular}{clc}
\toprule
\# & Topic & Date \\
\midrule
1 & Course Intro, FastAPI & 7/1 \\
2 & Basic Statistics \zhint{(统计)} & 14/1 \\
3--5 & Data Representation \zhint{(数据表示)} & 14--28/1 \\
6--7 & Data Modeling \zhint{(建模)} & 4--11/2 \\
8 & Data Visualization \zhint{(可视化)} & 25/2 \\
9 & Case: Business \zhint{(商业)} & 4--11/3 \\
10 & Case: Text Mining \zhint{(文本)} & 11--18/3 \\
11 & Case: Multimedia \zhint{(多媒体)} & 18--25/3 \\
12 & Case: Game \zhint{(游戏)} & 25/3--1/4 \\
13 & Summary \zhint{(总结)} & 1/4 \\
\# & Final Exam & 15/4 \\
\bottomrule
\end{tabular}
\end{center}

% ============================================================
\section{Python Environment \zhint{(Python环境)}}
% ============================================================

\subsection{What is Python? \zhint{(Python是什么)}}
\begin{itemize}[leftmargin=*, nosep, itemsep=0pt]
    \item High-level interpreted language (1991, Guido van Rossum)\\
    \zhint{高级解释型语言，代码简洁易读}
    \item Current version: Python 3.13 \zhint{(本课用Anaconda 3.13)}
\end{itemize}

\subsection{Hello World \zhint{(第一个程序)}}
\begin{lstlisting}
# In Python interpreter
>>> print("Hello World!")
Hello World!
\end{lstlisting}

\begin{lstlisting}
# As a script: helloworld.py
print("Hello World!")
# Run: $ python helloworld.py
\end{lstlisting}

\subsection{Common Uses \zhint{(Python能干啥)}}
\begin{itemize}[leftmargin=*, nosep, itemsep=0pt]
    \item Web backends \zhint{(网站后端)}: YouTube, Reddit
    \item Data science \& ML \zhint{(数据科学)}: PyTorch
    \item Scientific computing \zhint{(科学计算)}
    \item Data visualization \zhint{(数据可视化)}
\end{itemize}

\subsection{Setup \zhint{(推荐安装)}}
\begin{itemize}[leftmargin=*, nosep, itemsep=0pt]
    \item \textbf{Anaconda} --- 1500+ packages, easy to manage\\
    \zhint{自带大量数据科学包，装好就能用}
    \item \textbf{IDE}: VS Code \zhint{(免费，好用)}
\end{itemize}

% ============================================================
\section{Reference Books \zhint{(参考书)}}
% ============================================================

\begin{enumerate}[leftmargin=*, nosep, itemsep=2pt]
    \item \textit{Practical Statistics for Data Scientists}\\
    \zhint{统计学50个核心概念}
    \item \textit{Python for Data Analysis}\\
    \zhint{Pandas/NumPy实战数据处理}
    \item \textit{Intro to ML with Python}\\
    \zhint{Python机器学习入门}
    \item \textit{Data Science and Big Data Analytics}\\
    \zhint{数据科学全流程+大数据分析}
    \item \textit{Deep Learning} (Goodfellow)\\
    \zhint{深度学习经典教材（花书）}
\end{enumerate}

\subsection{Python Resources \zhint{(学习资源)}}
\begin{itemize}[leftmargin=*, nosep, itemsep=0pt]
    \item \texttt{docs.python.org/3/} \zhint{(官方文档)}
    \item \texttt{learnpython.org} \zhint{(交互式教程)}
    \item \texttt{w3schools.com/python} \zhint{(入门教程)}
\end{itemize}

\end{multicols}
\end{CJK}
\end{document}
